\documentclass[../../main]{subfiles}

\usepackage{tabularx}
\usepackage{caption}
\usepackage{eurosym}

\begin{document}

\title{Financial Transaction Tax}
\author{Gunther Capelle-Blancard}
\date{July 6, 2026}
\maketitle


A European financial transaction tax (FTT) is often presented as a politically attractive but technically fragile idea. This is misleading. The proposal considered here is not a utopian attempt to tax global finance in the abstract. It is a pragmatic equity transfer tax, modelled on devices that already exist and operate in major financial markets. The United Kingdom has applied a stamp duty on share transactions for more than three centuries. France introduced its current FTT in 2012, followed by Italy in 2013 and Spain in 2021; beyond Europe, comparable instruments also exist in major financial centres such as Hong Kong and Switzerland, and in around thirty countries worldwide. These examples show that a well-designed FTT can be implemented in open, sophisticated and highly internationalised financial markets without undermining their functioning.

%In the context of this report, the FTT is considered primarily as a fiscal instrument. The repayment of NextGenerationEU has created a need for credible new own resources that can strengthen the EU's long-term budgetary sustainability while remaining consistent with broader policy objectives. From this perspective, an equity transfer tax is an attractive candidate: it can raise revenue from a large, concentrated and relatively undertaxed tax base, while relying on mechanisms that are already in place in several major financial markets.

%This fiscal framing, however, should not reduce the FTT to a mere revenue device. The tax also has a strong theoretical and normative rationale. Since Keynes, Tobin and Stiglitz - to name a few -, one central argument has been that a small transaction cost can discourage excessive short-term trading without penalising long-term investment. The FTT should not be presented as a universal solution to financial instability or as a substitute for financial regulation, but it can act as a modest corrective instrument in markets where trading volumes have grown far beyond what is needed for ordinary investment and capital allocation.

%The FTT also has an attractive distributional profile compared with taxes on labour or broad-based consumption. Equity ownership and frequent trading are concentrated among wealthier households, institutional investors and financial intermediaries. It is therefore consistent with tax fairness, while also reflecting the broader costs that financial markets can impose on society through instability, excessive short-termism and recurrent public interventions in times of crisis. A well-designed FTT is theoretically grounded, non-distortive and inexpensive to collect.

The proposal is to introduce a harmonised EU tax on transfers of ownership of shares issued by companies incorporated in Member States, at a rate of 0.5\%, in line with the UK stamp duty.\footnote{The possibility that some new firms might incorporate outside the EU to avoid the tax cannot be entirely ruled out, but it is unlikely to be a major margin of avoidance. The tax would apply only to secondary-market transfers of shares, not to primary issuance, and its rate would be very low relative to the legal, fiscal, regulatory and governance costs associated with incorporating outside the firm's main economic area. For most firms, access to investors, legal certainty, listing conditions and corporate taxation are far more important than a modest tax on future share transactions.} Primary market issuance and genuine market-making activities could remain exempt, as in existing national schemes. The tax would be based on a robust legal and operational design: it would be attached to the security itself, rather than to the location of the broker, the trading venue or the investor. This is the key lesson from existing FTTs: such taxes can work, but only if they are carefully designed so as to limit avoidance and ensure enforceability.

The scope is deliberately narrow. At this stage, the tax would focus only on equities, not because derivatives, foreign-exchange transactions or crypto-assets should remain outside the scope of financial taxation, but because equity transfers provide the most robust and operationally tested base. This also draws a lesson from the earlier European Commission proposal. The 2011--2013 project was much broader: it covered shares, bonds and derivatives, and would have applied to gross transactions before netting and settlement. It was expected to raise much larger revenues, but it also triggered stronger technical and political objections and ultimately failed to secure agreement among Member States. The present proposal therefore does not seek to reopen the full debate on a comprehensive European FTT. It starts with the most feasible component, while leaving broader extensions (in particular the taxation of intraday transactions and other securities transactions) for a later stage.

Under conservative assumptions, the proposed EU equity transfer tax would raise approximately EUR 8 billion per year. The revenue estimate is deliberately prudent. It does not rely on speculative assumptions. Instead, it is anchored in the observed performance of the French FTT in 2024 and then applied to equity trading volumes in other Member States. This approach incorporates, by construction, the existing exemptions, leakages and the narrow effective tax base observed in France, including the exclusion of intraday transactions. 

In sum, an EU equity transfer tax would be theoretically grounded, technically feasible, fiscally useful and fair. It would not solve all budgetary challenges, and it should not be oversold. But as one component of a broader package of new own resources, it has a strong comparative advantage: it is based on a tested instrument, it targets a highly concentrated and undertaxed tax base, and even under very conservative assumptions it could generate around EUR 8 billion per year.

The following sections provide the detailed background supporting this proposal. It first presents the rationale and design of the equity transfer tax, before reviewing existing national policies and the earlier European FTT project. They then set out the revenue-estimation methodology and Member State-level results, responds to common objections, review public opinion, and finally discusses possible extensions, notably to intraday transactions.


%Political feasibility remains the main obstacle. A new EU own resource requires unanimity in the Council and approval by all Member States in accordance with their constitutional requirements. This is demanding, and the failure of the previous European FTT project shows that political resistance should not be underestimated. Yet the present proposal is more feasible precisely because it is narrower, more familiar and easier to implement. It does not require agreement on derivatives, foreign-exchange transactions or the taxation of all intraday market activity. It builds instead on instruments already used in several Member States.

%Finally, the proposal could serve as a platform for future extensions. The first and most natural extension would be to include intraday equity transactions. These transactions are largely outside the tax base, even though they include many of the most short-term strategies. Extending the tax to intraday trades would broaden the base, improve coherence and increase revenues without raising the statutory rate. The second extension would be to apply the FTT to other financial transactions.


%\newpage 

%\textbf{\LARGE{Appendix}}

\section{FTT: The rationale}

FTTs are neither new nor untested. They have existed for centuries and remain in force in many major financial centres. The most prominent example is the United Kingdom’s stamp duty on share transactions, introduced in the seventeenth century and still applied today to purchases of shares issued by UK companies. Similar taxes are also used in France, Italy, Spain and a number of other jurisdictions such as Switzerland, Singapore, Hong Kong, South Korea, etc. The relevant question is therefore not whether an FTT can exist, but how it should be designed \citep{CapelleBlancard2024}.

The proposal considered here is deliberately simple. It consists in introducing, at EU level, a harmonised tax on transfers of ownership of shares, modelled on existing national schemes such as the UK stamp duty and the French FTT. The tax would apply only to equity transactions, not to bonds, derivatives or foreign-exchange transactions. It would be levied at a rate of 0.5\%, in line with the UK stamp duty. The taxable event would be the transfer of ownership; in other words, the proposal does not require taxing all intraday transactions, but focuses on settled net acquisitions of shares. This proposal could serve as a platform for future extensions.

A key design feature is the issuance principle. Under this principle, the tax applies to transactions in shares issued by companies incorporated in the taxing jurisdiction, regardless of the residence or nationality of the investor and regardless of where the transaction takes place. This is the logic underpinning the UK stamp duty and the French FTT. It is crucial for enforceability: the tax is linked to the legal transfer of ownership of securities, not to the location of the trading platform or intermediary. This design considerably limits avoidance opportunities and explains why such taxes can be maintained even in highly internationalised financial markets \citep{CapelleBlancardPersaud2025}.

The economic rationale for an FTT is also well established. \citet{Keynes1936} argued that a tax on transactions could help reduce the excessive dominance of short-term speculation over enterprise. \citet{Tobin1978} later proposed a levy on currency transactions to ``throw sand in the wheels'' of international finance, while \citet{Stiglitz1989} defended the use of tax policy to curb speculative short-term trading. These arguments differ in scope, but they share a common intuition: when trading is excessively short-term, the private incentives of market participants may diverge from the social usefulness of financial markets. A small transaction cost can therefore serve both as a source of public revenue and as a modest corrective device.

For the purpose of new EU own resources, however, the fiscal argument is central. Existing FTTs are best understood not primarily as instruments of financial repression, but as low-rate taxes on a large and concentrated tax base. Their purpose is to generate revenue from financial market activity while preserving investment and market functioning. Evidence on the French FTT suggests that the tax reduced trading volumes moderately, by around 10--20\%, but without significant effects on market liquidity or volatility \citep{ColliardHoffmann2017}. This is important: the experience of recent European FTTs does not support the view that such taxes are incompatible with deep and efficient financial markets.

The proposal is therefore intentionally pragmatic. Rather than reviving the broader and more politically contentious idea of taxing all financial instruments, it builds on what already exists and works: a narrow, enforceable levy on equity transfers, based on the issuance principle, at a low nominal rate, with standard exemptions for primary issuance and market-making activities where appropriate. This approach would allow the EU to create a new own resource from a tried-and-tested tax instrument, while avoiding many of the objections usually directed at more comprehensive versions of the FTT.

\section{Existing national policies}

The proposed EU own resource would not be created in a vacuum. Several European countries already tax financial transactions, and the most relevant examples are not marginal cases: they include the United Kingdom, France, Italy and Spain. 

The United Kingdom remains the main benchmark. Its stamp duty on share transactions dates back to 1694 and is still in force. It applies at a rate of 0.5\% to purchases of shares issued by UK companies. The UK stamp duty raises \pounds{}4.32 billion in 2024--25, illustrating that such instruments can provide stable and non-negligible revenues. The key feature of the British model is the issuance principle: the tax applies because the share is issued by a UK company, irrespective of the nationality or residence of the investor and irrespective of where the transaction is executed. Primary market issuances, market-making activities and shares listed on the Alternative Investment Market are exempt. The UK experience is particularly important because it shows that a major financial centre can maintain a sizeable equity transaction tax over a very long period without undermining the development of its financial sector.

France introduced its current FTT in 2012. The French tax also follows the issuance principle. It applies to acquisitions of shares issued by companies whose registered office is in France and whose market capitalisation exceeds EUR 1 billion.\footnote{The list of companies subject to the tax is updated annually.} The tax is paid by the buyer and it is calculated on net acquisitions recorded at the end of the trading day. This means that the French FTT is not, strictly speaking, a tax on transactions, but a tax on transfers of ownership. Intraday transactions that are opened and closed within the same day are therefore effectively excluded from the tax base.\footnote{The French FTT formally has three components: a tax on share transfers, a tax on high-frequency trading, and a tax on naked sovereign credit default swaps. In practice, however, only the first component is effective. The high-frequency trading component is too narrow to generate meaningful revenue, while the naked CDS component became largely irrelevant after EU restrictions on such transactions.} The rate of the French FTT was initially 0.2\%, then increased to 0.3\% in 2017, and 0.4\% in 2025. The French case is especially useful for the present proposal because it shows how a modern equity FTT can be implemented unilaterally, even in an open financial market.

Italy introduced a similar tax in 2013, shortly after France. It applies to shares of Italian companies above a market-capitalisation threshold of EUR 500 million. The Italian design is broader in some respects, since it also includes a tax on derivatives and a specific component targeting high-frequency trading. It also differentiates between transactions executed on regulated markets and those executed elsewhere: the rate is 0.2\% for transactions on regulated markets and 0.4\% otherwise (since 2026). This feature is important because it can encourage trading to move towards more transparent and regulated venues. Italy therefore illustrates a possible refinement of the French model, even though the proposal considered here remains focused on a simple equity transfer tax.

Spain reintroduced an FTT in 2021, after having abolished an earlier tax in 1988. The Spanish tax follows the same broad family of instruments as the French and Italian taxes: it focuses on equity transactions and provides another recent example of a Member State implementing an FTT independently. Together, the French, Italian and Spanish cases show that unilateral implementation within the EU is feasible and that such taxes do not require a fully global agreement in order to operate.

Other European countries also have FTTs. Belgium has had a tax on stock-exchange transactions since 1913. Its design, however, differs from the British, French and Italian models. It is linked to the involvement of a Belgian financial institution or to Belgian-resident investors, and it appears easier to circumvent than systems based on the issuance principle. Ireland also applies a stamp-duty-type tax, introduced in 1937 on the British model, with a rate of 1\%. Existing work also identifies Finland, Poland, Cyprus and Malta among European countries with FTT-type instruments. These cases confirm that securities transaction taxes are not exceptional within Europe, although their design, scope and revenue performance vary substantially.\footnote{Several European countries previously levied similar taxes but later abolished them, notably the Netherlands in 1990, Germany and Sweden in 1991, Norway in 1993, Denmark in 1999 and Austria in 2000. These historical precedents should be interpreted with caution, as the abolished taxes often differed substantially in design from issuance-based equity transfer taxes such as the UK stamp duty or the French FTT.}

Overall, the international landscape is heterogeneous. Some instruments, such as the UK stamp duty, the French FTT or the Spanish FTT, are close to an issuance-based equity transfer tax. Others are broader securities transaction taxes, stamp duties or financial transaction taxes applying to different instruments or market segments. 

Table~\ref{tab:existing_ftt} provides an illustrative overview of selected existing taxes on equity and securities transactions, without attempting to provide an exhaustive or fully harmonised taxonomy.

The main lesson from existing national policies is clear. The success of an FTT depends less on the abstract principle of taxing financial transactions than on the details of its design. Taxes based only on the location of the intermediary or the trading venue are more vulnerable to avoidance. By contrast, taxes based on the issuance principle are more robust because the tax follows the security itself. The EU proposal examined here should therefore draw primarily on the UK, French, Italian and Spanish experiences, while avoiding the weaknesses of more easily circumvented systems.

\begin{table}[htbp]
\centering
\caption{Selected existing taxes on equity and securities transactions}
\label{tab:existing_ftt}
\footnotesize

\begin{tabularx}{\textwidth}{p{2.4cm} p{2.4cm} X}
\hline
Jurisdiction & Main rate & Main scope and design features \\
\hline

France & 0.4\% &
FTT on acquisitions of equity securities issued by French companies with market capitalisation above \euro{}1 billion. Primary issuance and several market operations are exempt; intraday transactions are outside the effective tax base. \\

Italy & 0.2--0.4\% &
FTT on transfers of shares and equity-like instruments issued by Italian companies with market capitalisation above \euro{}0.5 billion. The lower rate applies to transactions executed on regulated markets or multilateral trading facilities. The Italian FTT also includes specific, albeit marginal, components for derivatives and high-frequency trading. \\

Spain & 0.2\% &
FTT on onerous acquisitions of shares in Spanish companies with market capitalisation above \euro{}1 billion. \\

Belgium & 0.12\%-1.32\% &
Tax on stock-exchange transactions. The applicable rate and cap depend on the type of financial instrument. This design is less close to a pure issuance-based equity transfer tax. \\

Ireland & 1\% &
Stamp duty on instruments or transfer orders transferring shares, stocks or marketable securities. \\

Greece & 0.1\% &
Tax on sales of shares listed on a Greek regulated market or multilateral trading facility. The rate was reduced from 0.2\% to 0.1\% from 2024. \\

\hline 

United Kingdom & 0.5\% &
Stamp Duty / SDRT on purchases or transfers of shares in UK companies. Electronic transactions are generally collected through CREST. This remains the main benchmark for an issuance-based equity transfer tax. \\

Switzerland & 0.15\% / 0.30\% &
Securities transfer stamp tax on purchases and sales of Swiss or foreign securities when a Swiss securities dealer acts as contracting party or intermediary. The lower rate applies to Swiss securities and the higher rate to foreign securities. \\

\hline 

Hong Kong & 0.1\%+0.1\% &
Stamp duty on the sale or purchase of Hong Kong stock. The duty is charged on every bought note and every sold note. \\

Singapore & 0.2\% &
Share stamp duty on transfers of shares, computed on the purchase price or value of the shares transferred. \\

South Korea & 0.15\%-0.35\% &
Securities transaction tax on transfers of Korean shares. Rates differ across listed and unlisted markets and have changed repeatedly in recent years. \\

China & 0.05\% &
Stamp duty on stock transactions, reduced from 0.1\% to 0.05\% from 28 August 2023. \\

Taiwan & 0.3\% &
STT levied on sellers of securities. Reduced rates or suspensions apply to some instruments and transactions. \\

India & 0.1\% &
STT applying to exchange-traded equity and derivatives transactions. Rates differ across delivery-based equity trades, intraday trades, futures and options. \\

%Malaysia & 0.1\% effective rate, capped &
%Stamp duty on contract notes for listed shares or stocks, with the effective duty reduced to 0.1\% and capped at RM 1,000 under the current remission order. \\

South Africa & 0.25\% &
Securities transfer tax on the transfer of listed and unlisted securities. \\

 \\

\hline
\end{tabularx}

\caption*{\footnotesize Notes: This table is illustrative rather than exhaustive. It focuses on existing taxes that are closest to an equity transaction or equity transfer tax. Rates and scope are simplified; exemptions, collection mechanisms and anti-avoidance rules differ substantially across jurisdictions.}

\end{table}


\newpage

\section{The (former) European FTT project}

The idea of a European financial transaction tax gained renewed political momentum after the 2007--2008 financial crisis. The crisis strengthened the view that the financial sector should make a fairer contribution to public finances and that new fiscal instruments could also help discourage socially unproductive short-term trading. The European Parliament called on the Commission to examine the issue in 2010, and the initiative was then strongly supported by France and Germany. In September 2011, the European Commission adopted a proposal for a Council directive establishing a common system of financial transaction tax across the European Union.

The Commission proposal was deliberately ambitious \citep{Gabor2016}. Its objective was not merely to create a new source of revenue, but also to harmonise existing national approaches, prevent fragmentation of the Single Market, ensure a fair contribution from the financial sector, and discourage transactions considered not to contribute to the economy. The proposal was based on a broad tax base and low rates. It would have applied to secondary-market transactions involving financial institutions, with minimum rates of 0.1\% for shares and bonds and 0.01\% for derivatives. The Commission described the approach as applying to all markets, all instruments and all actors. In particular, the proposed tax would have applied to gross transactions before netting and settlement, thereby including intraday transactions.

The expected revenue was substantial. For the EU as a whole, the Commission estimated that the proposal could raise around EUR 57 billion per year. Other estimates for the eleven Member States later involved in the enhanced cooperation procedure suggested annual revenues of around EUR 30--35 billion. These figures reflected the breadth of the initial proposal: a large share of expected revenue came not from the taxation of equity transactions alone, but from the inclusion of bonds and, especially, derivatives.

The project immediately encountered strong opposition. Because taxation requires unanimity in the Council, the proposal could not be adopted at EU-wide level. Several Member States were reluctant or hostile, especially those concerned about the impact on financial centres, market liquidity or the relocation of trading activity. The United Kingdom, Sweden and the Netherlands were among the most vocal opponents in the political debate. According to a contemporary review by the Fondation Robert Schuman, six countries eventually indicated formal opposition to the project in the context of the enhanced-cooperation decision: Bulgaria, Luxembourg, Malta, the United Kingdom, Sweden and the Czech Republic. Ireland and Latvia expressed reservations without formally opposing the principle.

As unanimity proved impossible, the debate shifted towards enhanced cooperation. In 2012, eleven Member States requested authorisation to proceed together: Belgium, Germany, Estonia, Greece, Spain, France, Italy, Austria, Portugal, Slovenia and Slovakia. The Council authorised enhanced cooperation in January 2013, and the Commission tabled a revised proposal in February 2013 for these participating countries. Estonia later withdrew from the process, leaving ten participating Member States. The use of enhanced cooperation was politically significant, as it would have been the first case of such a procedure in the field of taxation.

Despite this narrower framework, negotiations never reached a conclusion. The participating Member States continued to disagree on the scope of the tax, the treatment of derivatives, the residence and issuance principles, the allocation of revenues and the possible effects on non-participating Member States. The proposal remained blocked in the Council for years. In 2023, the Commission stated that the prospects of reaching an agreement were limited and that there was little expectation of an agreement in the short term. In its 2026 work programme, the Commission indicated its intention to withdraw the proposal.

The proposal examined here is much less ambitious than the earlier European project. It does not seek to tax all financial instruments, all markets and all financial actors. It does not include bonds, derivatives or foreign-exchange transactions. It also does not attempt to tax all gross intraday transactions. Instead, it focuses on equity transfers, following the logic of existing national schemes such as the British stamp duty and the French FTT. The rate considered here, 0.5\%, is higher than the 0.1\% rate proposed by the Commission for shares and bonds, but the tax base is much narrower. The proposal is therefore best understood as a pragmatic equity transfer tax rather than as a comprehensive European FTT.

This narrower design has two advantages. First, it relies on mechanisms that have already been implemented in major financial markets. Second, it avoids reopening all the technical and political controversies that contributed to the failure of the earlier European project. The objective is not to revive the full 2011--2013 proposal, but to draw a practical lesson from it: an overly broad FTT may be politically fragile, whereas a simpler tax on equity transfers, based on the issuance principle and modelled on existing national systems, is more likely to be feasible.

\section{Revenue potential at Member State level}

The revenue estimates presented here follow an inference-based approach, rather than a purely hypothetical one \citep{CapelleBlancardPersaud2025}. The method is anchored in the observed revenue of the French FTT and then applied to equity trading volumes in other Member States. This is deliberately conservative: rather than assuming that the full statutory tax rate applies to all market transactions, it starts from what the French tax actually raises under the existing design, with its exemptions, its daily netting mechanism and the exclusion of intraday trading.

The reference year is 2024. In France, the FTT raised EUR 1,860 million, including EUR 528 million allocated to the Agence française de développement. According to ESMA\footnote{\textit{ESMA statistics on securities and markets}, ESMA11-239717167-20877.}, the total value of equity transactions in France in 2024 was EUR 3,502,800 million. The implicit tax ratio is therefore:

\[
\frac{1,860}{3,502,800} = 0.0531\%.
\]

This figure is much lower than the statutory rate of the French FTT, which was 0.3\% in 2024 (0.4\% since April 2025). It implies that the effectively taxed base represents only:

\[
\frac{0.0531\%}{0.3\%} = 17.7\%
\]

\noindent of total equity trading volumes. In other words, the implicit exemption or out-of-scope rate is approximately 82.3\%. This reflects the narrowness of the current French design: the tax applies to transfers of ownership recorded at the end of the day, not to all intraday transactions. Paradoxically, the most short-term transactions are therefore largely left outside the tax base.

The same implicit tax ratio is then applied to the equity trading volume of each Member State.\footnote{For some Member States, ESMA reports no equity trading volume in the dataset used for this exercise. These entries are therefore kept at zero for consistency, but they should not be interpreted as implying the complete absence of equity market activity.} The estimate at a 0.3\% statutory rate is given by:

\[
R_i^{0.3} = T_i \times \frac{1,860}{3,502,800},
\]

\noindent where \(T_i\) is the total equity trading volume in Member State \(i\). The estimate at a 0.5\% rate is obtained by assuming a proportional increase in revenue:

\[
R_i^{0.5} = R_i^{0.3} \times \frac{0.5}{0.3}.
\]

Table~\ref{tab:revenues_member_states} reports the revenues estimated using this approach for each EU Member State.\footnote{The Member State allocation is based on observed equity trading volumes reported in the ESMA dataset, rather than on the headquarters of the issuing firms. It should therefore be read as an operational approximation, not as a precise legal allocation under the issuance principle. Existing national FTTs are not deducted from these gross estimates.} Overall, an EU-wide equity transfer tax at a 0.5\% rate would raise approximately EUR 8 billion per year. This should be interpreted as an estimate for a narrow tax on equity transfers, not as an estimate for a broader FTT including bonds, derivatives or intraday transactions.\footnote{These illustrative figures differ from previous estimates of FTT revenues \citep{CapelleBlancardPersaud2025} for methodological reasons. In earlier work, transaction volumes were based on a combination of World Bank and Refinitiv data for 2022, and revenues were estimated by assuming that 15\% of total turnover would remain taxable after exemptions and that trading volumes would fall by 20\% after the introduction of the tax. At a 0.5\% statutory rate, this amounts to applying an effective revenue ratio of approximately 0.06\% to total trading volumes. The estimates reported here use ESMA (and WFE) trading volumes instead, and apply the implicit French revenue ratio observed in 2024, scaled from the French statutory rate to a 0.5\% rate. The difference between the two sets of results therefore reflects both the use of a different data source for trading volumes and a different effective tax ratio. These figures should be read as illustrative orders of magnitude rather than as directly comparable estimates.}

The distribution of revenues would be highly concentrated. France alone would account for about 37\% of total EU revenue, followed by the Netherlands with 23.5\% and Germany with 19.1\%. Together, these three countries would generate almost 80\% of the total. Including Italy brings the share to around 86\%. This concentration reflects the geography of equity trading within the EU rather than any specific assumption about tax capacity or political preferences.\footnote{As an illustrative benchmark, applying the same implicit revenue ratio to WFE equity trading volumes gives much larger orders of magnitude for major non-EU markets: approximately EUR 116 billion for the United States, EUR 6.9 billion for Japan, EUR 28.6 billion for China, and EUR 178 billion at the world level. These figures should be interpreted with great caution. WFE trading volumes are not directly comparable with the ESMA data used for the EU Member State estimates. They are therefore provided only as indicative orders of magnitude, not as directly comparable revenue estimates.}

This concentration has two implications. First, the fiscal potential of the measure depends on a small number of large and liquid equity markets. Second, from an own-resources perspective, the tax would not mechanically generate revenue in proportion to Member States' GDP or population. This is not necessarily a weakness: the tax is designed to capture activity in financial markets, and market activity itself is highly concentrated. But it does mean that the allocation of revenues and the political presentation of the measure would need to be carefully designed.

\newpage

\begin{longtable}{lrrr}
\caption{Estimated revenues by EU Member State}
\label{tab:revenues_member_states}\\
\hline
Member State & Code & Equity trading volume & Estimated revenue \\
             &      & EUR million & EUR million, 0.5\% rate \\
\hline
Austria & AT & 29,800 & 26.4 \\
Belgium & BE & 69,000 & 61.1 \\
Bulgaria & BG & 200 & 0.2 \\
Cyprus & CY & 900 & 0.8 \\
Czechia & CZ & 4,300 & 3.8 \\
Germany & DE & 1,802,500 & 1,595.2 \\
Denmark & DK & 128,600 & 113.8 \\
Estonia & EE & 100 & 0.1 \\
Spain & ES & 271,900 & 240.6 \\
Finland & FI & 91,900 & 81.3 \\
France & FR & 3,502,800 & 3,100.0 \\
Greece & GR & 25,700 & 22.7 \\
Croatia & HR & 300 & 0.3 \\
Hungary & HU & 6,700 & 5.9 \\
Ireland & IE & 272,800 & 241.4 \\
Italy & IT & 615,200 & 544.5 \\
Lithuania & LT & 100 & 0.1 \\
Luxembourg & LU & 0 & 0.0 \\
Latvia & LV & 0 & 0.0 \\
Malta & MT & 0 & 0.0 \\
Netherlands & NL & 2,215,800 & 1,961.0 \\
Poland & PL & 71,200 & 63.0 \\
Portugal & PT & 28,100 & 24.9 \\
Romania & RO & 2,200 & 1.9 \\
Sweden & SE & 277,000 & 245.1 \\
Slovenia & SI & 500 & 0.4 \\
Slovakia & SK & 0 & 0.0 \\
\hline
European Union & EU27 & 9,417,600 & 8,334.6 \\
\hline
\end{longtable}

\newpage

\section{Responses to common objections}

The debate on FTTs is marked by a recurring set of objections. Many of them are serious and should not be dismissed. However, they often refer to very broad FTT schemes, whereas the proposal considered here is narrower: a tax on equity ownership transfers, based on the issuance principle, at a moderate rate. This distinction is crucial. The question is not whether any possible FTT would be desirable, but whether this specific design can provide a robust and workable own resource.


\begin{enumerate}

    \item \textbf{The FTT would reduce market liquidity.}

    This objection is partly correct but often overstated. A FTT increases the cost of trading and is therefore expected to reduce trading volumes \citep{IMF2010, ColliardHoffmann2017}. This is not necessarily a problem. A fall in very short-term trading is one of the intended effects of such a tax, especially when a large share of market activity consists of rapid, low-margin transactions whose private profitability may exceed their social value. Some high-frequency strategies may improve liquidity and price discovery, but others mainly reflect a technological arms race for speed, latency arbitrage, order anticipation or very short-lived quote revisions. These activities can impose costs on other market participants, increase the resources devoted to trading infrastructure, and generate little additional information for long-term capital allocation.\footnote{The literature on individual investors  shows that more frequent trading often reduces net performance \citep{DeBondtThaler1995, Shleifer2000,Odean1999,BarberOdean2000, BarberOdean2001}, while recent work on high-frequency trading emphasises the private gains but limited social value of ever-faster execution \citep{BudishCramtonShim2015, AquilinaBudishONeill2022}.}

    The relevant question is not whether turnover declines, but whether useful liquidity is harmed. Existing evidence on modern, low-rate equity FTTs suggests that turnover may fall, but without clear evidence of significant deterioration in liquidity or volatility \citep{ColliardHoffmann2017}. 
    
    The Swedish case is often used as an argument against FTTs, but its main problem was poor design: the tax was linked to domestic intermediaries and was therefore easy to avoid by relocating transactions. It is not evidence that a well-designed issuance-based tax necessarily decreases liquidity. The proposal considered here focuses on ownership transfers and can preserve exemptions for genuine market-making activities.

    \item \textbf{The FTT would increase volatility.}

    The theoretical effect of an FTT on volatility is ambiguous. If all market participants are perfectly rational and liquidity is always beneficial, then a tax may appear harmful. But if markets also include noise trading, herding, excessive short-termism or self-reinforcing speculation, then reducing some transactions may improve rather than worsen price formation \citep{Matheson2011, McCullochPacillo2011}.

    Empirically, the claim that FTTs mechanically increase volatility is not supported by the experience of existing low-rate equity taxes. 

    \item \textbf{The FTT would raise the cost of capital for firms.}

This objection relies on a plausible but indirect mechanism: a transaction tax could reduce trading volumes; lower trading volumes could reduce market liquidity; lower liquidity could increase the return required by investors; and this, in turn, could raise firms' cost of capital. The first link in this chain is credible: an FTT is expected to reduce trading volumes. The key question, however, is whether this reduction in turnover translates into a meaningful deterioration in liquidity and then into a higher cost of capital.

The theoretical link between liquidity and expected returns is well established \citep{AmihudMendelson1986,Amihud2002}. But the empirical relevance of this channel for a low-rate equity transfer tax is much less clear. Evidence on the French FTT shows that trading volumes declined, but the effects on market quality were limited and did not provide evidence of a large disruption to liquidity. If the liquidity channel is weak or breaks at this stage, there is little reason to expect a material effect on firms' cost of capital.

Direct evidence is also reassuring. Fraichot (2017) studies the effect of an FTT on corporate cost of capital and finds an insignificant impact in highly liquid equity option markets. More recently, Do (2025) studies the 2012 French FTT and finds a positive effect on corporate investment, especially where the tax induced a shift from short-term to longer-term ownership. This suggests that any increase in transaction costs may be offset by a reduction in short-termism, rather than mechanically raising firms' financing costs \citep{Fraichot2017,Do2025}.

    \item \textbf{The FTT would harm investment and competitiveness.}

    This objection assumes that the presence of a FTT makes a financial centre unattractive. Existing experience suggests otherwise. The United Kingdom has maintained a stamp duty on share transactions for centuries while remaining one of the world's leading financial centres. France, Italy and Spain have also introduced modern equity FTTs without undermining the functioning of their equity markets.

    Competitiveness concerns should therefore be assessed in relation to the design of the tax. A badly designed tax can create distortions. A tax based on the issuance principle, applied at a moderate rate and limited to equity transfers, is much less likely to damage competitiveness than a tax based only on the location of intermediaries or trading venues.

    \item \textbf{The FTT would lead to relocation.}

    Relocation is a serious risk when the tax is based on the location of the broker, the exchange or the intermediary. This was the central weakness of the Swedish experiment in the 1980s. It is not the logic of the proposal considered here.

    Under the issuance principle, the tax follows the security, not the marketplace. A share issued by a company in a participating Member State remains within the scope of the tax even if it is traded through a foreign intermediary or on an alternative platform. This substantially limits the scope for relocation, because moving the transaction does not remove the underlying security from the tax base.

    \item \textbf{The FTT would be easy to avoid.}

    Avoidance is possible for any tax, but the risk is much lower when the tax is attached to the transfer of legal ownership. The British stamp duty illustrates this logic: the tax is connected to the recognition and settlement of ownership. The French system also shows that an equity transfer tax can be implemented in a modern, highly internationalised market.

    The main remaining issue is substitution towards untaxed instruments or arrangements, including derivatives. This is a valid concern. However, the proposal examined here is deliberately conservative and focuses on the most robust part of the tax base: equity ownership transfers. Additional anti-avoidance rules could be considered if substitution became significant, but the existence of possible avoidance margins is not a reason to abandon the tax altogether.

    \item \textbf{The FTT would hit ordinary savers.}

    This objection is often politically powerful but economically incomplete. The direct payer may be the buyer of the share, but the economic incidence of the tax depends on trading behaviour, portfolio structure and market intermediation. Ordinary savers typically trade much less frequently than professional intermediaries or short-term investors. For buy-and-hold investors, the effective annual cost of a small transaction tax is not significant.

The French case suggests that a substantial share of the tax would also be collected from foreign investors. In 2024, non-residents held 50\% of the capital of French-resident CAC 40 companies. Under the (plausible) assumption that trading shares are broadly proportional to ownership shares, roughly half of the revenue from an issuance-based tax on French shares would therefore be collected from non-resident investors. For an EU-wide tax, the share paid by non-EU investors would require a dedicated calculation, since intra-EU cross-border holdings would no longer count as external to the Union. The French evidence nevertheless shows that such taxes are not borne only by domestic households.

The redistributive profile of the tax is also important. Equity ownership is highly concentrated among wealthy households. In the euro area, the wealthiest 10\% of households own about 82\% of directly held listed equities; in France, the concentration is even higher, at close to 88\%.\footnote{https://www.afg.asso.fr/app/uploads/2025/02/Etude-AFG-OEE-2025.pdf} By contrast, direct share ownership is limited among the rest of the population. Even when the FTT is ultimately borne by households, it therefore falls disproportionately on wealthier investors and more frequent traders, rather than on ordinary savers. This makes the FTT more progressive than broad-based taxes on labour income or consumption.

    \item \textbf{The FTT would penalise market-making and useful liquidity provision.}

   This is a design issue rather than a decisive objection. The purpose of the tax is not to penalise genuine liquidity provision, but to tax equity ownership transfers while preserving the market functions that are necessary for orderly trading. Existing FTTs already address this issue through clear exemptions. The French FTT, for example, exempts market-making activities carried out by investment firms or credit institutions, including when these activities are conducted abroad. It also exempts market-making activities carried out on behalf of the issuer in order to support the liquidity of its securities. These exemptions are explicitly designed to protect useful liquidity provision.

The key challenge is therefore to calibrate exemptions carefully. If they are too narrow, useful liquidity provision may be affected. If they are too broad, the tax base is unnecessarily eroded. 


    \item \textbf{The FTT would be administratively complex to collect.}

    The practical experience of existing taxes suggests the opposite. Equity transfer taxes can be collected through settlement, reporting or central securities infrastructure. The UK stamp duty and the French FTT show that such taxes are technically feasible.

    That said, the quality of the collection mechanism matters. The French experience also shows that transparency, access to data and administrative control are important. A European own resource should therefore be designed with clear reporting obligations and strong cooperation between tax authorities, market supervisors and settlement infrastructures.

    %\item \textbf{The FTT would generate uncertain or disappointing revenues.}

    %Revenue uncertainty exists for all taxes, especially those linked to market activity. But the estimation method used here is deliberately conservative. It does not assume that all equity trading volumes are taxed. Instead, it starts from observed French revenue and applies the corresponding implicit tax ratio to other Member States.

    %This approach already incorporates the main exemptions and the fact that intraday transactions are largely outside the tax base. The resulting estimates should therefore be understood as estimates for a narrow equity transfer tax, not for a broad tax on all financial transactions. The proposal does not rely on optimistic assumptions about taxing the entire volume of market activity.

    %\item \textbf{The FTT would require global coordination.}

    %A global agreement would increase the revenue potential of an FTT, but it is not a precondition for implementation. The UK, France, Italy and Spain have all implemented equity transaction or transfer taxes without waiting for a global agreement.

    %The reason is again the issuance principle. If the tax is attached to shares issued in the taxing jurisdiction, it can operate even when investors, intermediaries or trading venues are located abroad. International coordination is desirable, but it is not technically indispensable for a narrow equity transfer tax.

    %\item \textbf{The FTT is the wrong instrument for financing EU priorities.}

    %The FTT should not be presented as a universal solution. It will not, by itself, repay NextGenerationEU, solve fiscal sustainability issues or regulate financial markets. But this is not the relevant benchmark. The question is whether it can be one own resource among several.

    %As one component of a broader package, an equity transfer tax has several advantages: it is already tested, it can raise revenue from a large and highly internationalised tax base, it has a redistributive profile, and it is connected to an activity that is currently only lightly taxed relative to its scale. It is therefore a plausible and coherent own-resource candidate.

    %\item \textbf{The FTT is politically unrealistic.}

    %The failure of the previous European FTT project is often cited as proof that the idea is politically impossible. This conclusion is too strong. The earlier EU proposal was much broader: it covered several categories of financial instruments, including derivatives, and raised complex questions about scope, allocation and cross-border effects.

    %The proposal examined here is much less ambitious. It is closer to existing national policies and focuses only on equity ownership transfers. This narrower design does not eliminate political difficulties, but it changes their nature. It is easier to defend a pragmatic extension of instruments that already exist in several major financial markets than to reopen the entire debate on a comprehensive European FTT.

\end{enumerate}


%\section{Assessment of political feasibility and EU decision-making procedures}

%The political feasibility of an EU equity transfer tax should be assessed with some caution. The measure is technically feasible and already exists in several European countries, but taxation remains one of the most politically sensitive areas of EU decision-making. The relevant question is therefore not whether such a tax can be designed or collected, but whether Member States can agree to turn it into a European own resource.

%From a legal and institutional point of view, two issues must be distinguished. First, if the tax is introduced as a genuine EU own resource, it falls under the own-resources procedure. Under Article 311 TFEU, the Council must act unanimously after consulting the European Parliament, and the decision does not enter into force until it has been approved by all Member States in accordance with their constitutional requirements.\footnote{Article 311 of the Treaty on the Functioning of the European Union.} This makes the procedure demanding, since a single Member State can block the adoption of a new own resource.

%Second, if the proposal requires harmonised rules on the taxation of financial transactions, it would also raise issues of tax harmonisation. The previous European FTT proposal was based on Article 113 TFEU, which concerns the harmonisation of indirect taxes necessary for the functioning of the internal market and the avoidance of distortions of competition. In that framework too, unanimity in the Council is required. The enhanced-cooperation route can reduce the number of participating countries, but it does not remove the need for agreement among those participating Member States. In the case of the 2013 FTT proposal, adoption of the implementing directive required unanimous agreement of the participating countries within the Council, after consultation of the European Parliament and the European Economic and Social Committee.\footnote{European Parliament, Legislative Train Schedule, \textit{Financial Transaction Tax}.}

%The historical record is therefore mixed. On the one hand, there has been clear political support for an FTT in several large Member States. France, Italy and Spain have already introduced national equity transaction taxes. Germany, France, Italy, Spain, Belgium, Austria, Greece, Portugal, Slovenia and Slovakia were part of the enhanced-cooperation process after the failure of the EU-wide proposal, with Estonia initially involved before withdrawing. On the other hand, the previous European project never reached agreement. The 2011 EU-wide proposal failed because unanimity could not be obtained, and the 2013 enhanced-cooperation proposal remained blocked for years. In 2023, the Commission stated that the prospects of reaching agreement were limited, and its 2026 work programme indicated an intention to withdraw the pending FTT proposal.\footnote{European Parliament, Legislative Train Schedule, \textit{Financial Transaction Tax}; European Commission, Work Programme 2026.}

%Past national positions also suggest where political resistance may arise. During the earlier debate, Bulgaria, Luxembourg, Malta, Sweden, the Czech Republic and the United Kingdom formally opposed the move towards enhanced cooperation, while Ireland and Latvia expressed reservations without opposing the principle.\footnote{Fondation Robert Schuman, \textit{The European Tax on Financial Transactions: from principle to implementation}, 2013.} The United Kingdom is no longer a Member State, but the other cases remain relevant as indicators of likely concerns: protection of domestic financial centres, fear of relocation, reluctance towards new EU-level taxation, and scepticism about partial implementation.

%The proposal examined in this annex is, however, significantly less ambitious than the earlier European FTT. It does not cover all financial instruments, all markets or derivatives. It does not seek to tax all gross intraday transactions. It is instead limited to equity ownership transfers, modelled on existing national schemes and based on the issuance principle. This narrower design improves political feasibility in three ways. First, it is easier to explain: it is not a new and untested tax, but an extension and harmonisation of instruments already used in major markets. Second, it reduces the scope of technical objections, especially those related to derivatives and market-wide relocation. Third, it produces more conservative revenue estimates, because it does not rely on taxing the full volume of financial transactions.

%There remains, nevertheless, an important political difficulty. The estimated revenues are highly concentrated in a small number of Member States, especially France, the Netherlands and Germany. This concentration reflects the geography of equity trading in Europe, but it may complicate negotiations if the tax is perceived as falling disproportionately on a few financial centres while the revenues accrue to the EU budget as a whole. This issue is not specific to the FTT, but it would need to be addressed explicitly in the political presentation of the proposal.

%Overall, the political feasibility of the proposal can be described as moderate rather than high. An EU-wide own resource based on an equity transfer tax would face the same institutional obstacle as all new own resources: unanimity and national approval. The failure of the previous European FTT also shows that political resistance should not be underestimated. However, the narrower design proposed here makes the measure more feasible than the earlier Commission project. It is closer to existing national experience, less exposed to technical controversy, and better suited to being presented as a pragmatic own-resource option rather than as a comprehensive reform of financial taxation.

%A realistic political strategy would therefore avoid reopening the full 2011--2013 FTT debate. It should instead emphasise the limited scope of the proposal, the precedent of existing national equity taxes, the robustness of the issuance principle, and the specific need to create credible new own resources for the repayment of NextGenerationEU.

\section{Review of existing surveys on public opinion}

Public opinion is one of the relative strengths of the FTT proposal. The available evidence should nevertheless be interpreted with caution. The most detailed EU-wide survey is now somewhat dated, while the most recent survey evidence covers only a limited number of Member States. Taken together, however, the results point to a consistent pattern: financial transaction taxes tend to attract broad public support, especially when they are presented as a way to make the financial sector contribute to public goods, development, climate action or the costs of crises.

The main EU-wide reference remains the 2011 Eurobarometer survey commissioned by the European Parliament in the aftermath of the financial crisis.\footnote{European Parliament, \textit{Eurobarometer 75.2: Europeans and the Crisis}, Analytical synthesis, June 2011.} In that survey, 61\% of Europeans supported the principle of a tax on financial transactions, while 26\% were opposed and 13\% expressed no opinion. Support was higher in euro area countries, at 63\%, than in non-euro area countries, at 54\%. The wording of the question explicitly stated that the tax would be very low and would apply to transactions between financial operators, not to the general public.

The same survey also provides useful information on the reasons behind public support. Among respondents in favour of the principle of the tax, the main motivation was to combat excessive speculation and help prevent future crises, followed by the idea that financial actors should contribute to the costs of the crisis. Reducing public deficits and financing innovative policies, including climate, environment and development aid, were also mentioned, but less frequently. This suggests that the public appeal of the FTT is linked not only to revenue-raising, but also to a broader perception of fairness in the taxation of the financial sector.

The objections expressed by opponents in the 2011 survey were also informative. Among those opposed to introducing the tax at EU level as a first step, the main concern was feasibility: some respondents believed that such a tax could only be introduced globally. Other concerns related to competitiveness, the risk that only European financial actors would contribute, and possible capital outflows from the EU. These objections closely mirror the arguments raised in the policy debate. They also reinforce the importance of explaining why the proposal considered here is narrower, issuance-based, and modelled on existing national systems.

More recent evidence is provided by a 2024 survey conducted by Focus 2030 and Stack Data Strategy in France, Germany and Italy, ahead of the European Parliament elections.\footnote{Focus 2030 and Stack Data Strategy, \textit{Survey: citizen opinions in France, Italy and Germany on the role of the European Union in international development ahead of the elections}, March 2024.} The survey was carried out between 16 and 22 February 2024, with representative samples of around 1,500 respondents in each of the three countries. It found that 57\% of respondents supported a European financial transaction tax, while 80\% were either in favour or would not oppose such a tax if it were used to finance the fight against global poverty and climate change. Only 13\% opposed the measure. Support, or non-opposition, was highest in France, at 82\%, with only 9\% opposed.

%This recent survey is not a substitute for a full EU27 Eurobarometer, but it is politically relevant. France, Germany and Italy are among the largest Member States and were central to previous European discussions on the FTT. The 2024 results suggest that, in these countries at least, a European FTT linked to global public goods and climate finance does not appear to carry a major public-opinion cost. The measure remains more popular among left-leaning voters, but opposition remains limited overall.

Available surveys suggest that the FTT is more popular than many other tax instruments, especially when framed as a modest levy on financial transactions rather than a tax on ordinary households. From a political communication perspective, the strongest framing is therefore not technical harmonisation, but fairness: a small contribution from financial-market activity to finance common European and global priorities.


\section{Possible extensions}

The proposal examined here is intentionally limited to equity ownership transfers. This narrow scope is useful for adoption, but it could later serve as a platform for two main extensions: first, the inclusion of intraday equity transactions; second, the extension of the tax base to other financial instruments, such as derivatives or crypto-assets.\footnote{Another possible extension would be a green FTT, either through the allocation of revenues to climate and environmental objectives or through differentiated rates linked to issuers' environmental performance; see \citep{CapelleBlancardGreenFTT}.}

\subsection{Extension to intraday equity transactions}

The most natural extension would be to include intraday equity transactions. Under the French model, the tax applies to net acquisitions recorded at the end of the trading day. Transactions opened and closed within the same day are therefore outside the effective tax base, including many high-frequency and very short-term strategies. This is a major limitation: the tax applies to investors who hold shares at least until the end of the day, while many of the shortest-term transactions remain untaxed.

Including intraday transactions would broaden the base, increase revenues and improve the coherence of the tax. In the French case, the effectively taxed base represents only a small share of total equity trading volumes, precisely because the tax is collected on transfers of ownership rather than on all transactions. From an economic point of view, it is difficult to justify taxing longer-term investors while leaving the fastest and most speculative transactions outside the scope of the tax.

The main obstacle is operational rather than conceptual \citep{CapelleBlancard2024FrenchFTT}. A tax on end-of-day ownership transfers can be collected through central securities depositories such as Euroclear. A tax on intraday transactions, by contrast, requires reliable information at the level of order execution and trading venues. This would imply a change in the collection architecture. One possible route would be to rely more directly on transaction-level reporting under MiFID II, which already provides regulators with detailed information on market transactions.

This extension should not be confused with the taxation of genuine liquidity provision. Market-making is an economic function, not a trading horizon. Some intraday transactions contribute to liquidity, but many others do not. Existing FTTs already include targeted exemptions for market-making activities, including in France. A blanket exclusion of all intraday trading is therefore not the right way to protect liquidity. It amounts to adding a broad and poorly targeted exemption on top of already existing protections. A well-designed EU FTT should preserve clear market-making exemptions, but should not exclude all intraday trading by default.

\subsection{Extension to other financial instruments}

A second extension would consist in broadening the tax base beyond equity transfers. The earlier European Commission proposal covered not only shares, but also bonds and derivatives. Such a broader base would raise larger revenues and reduce substitution towards untaxed instruments. It would also address a legitimate concern: if only equity transfers are taxed, investors may replicate some exposures through derivatives or other instruments that are economically close to the taxed transaction.

This extension is, however, more complex than the inclusion of intraday equity transactions. Derivatives raise difficult questions of valuation, netting, maturity, notional amounts and economic equivalence. A simple tax on the notional value of derivatives would not be appropriate for all contracts, while a tax on premia or market values may be harder to harmonise and enforce. Crypto-assets raise additional issues, including the identification of taxable events, the location of platforms and intermediaries, and the treatment of decentralised transactions.

For these reasons, the core proposal should remain focused on equity transfers. Once the legal and administrative framework is established, the EU could consider targeted extensions to instruments that are close substitutes for equity transactions, especially equity derivatives. Broader extensions to bonds, foreign exchange or crypto-assets should be treated as separate policy choices, requiring specific design and impact assessment. The priority is therefore sequential: first establish a robust equity transfer tax; then broaden the base where substitution risks, feasibility and revenue potential justify doing so.



\bibliographystyle{apalike}
\bibliography{FTT_references}


\end{document}